% Part: first-order-logic
% Chapter: axiomatic-proofs
% Section: identity

\documentclass[../../../include/open-logic-section]{subfiles}

\begin{document}

\olfileid[hi]{fol}{axd}{ide}

\olsection{\usetoken{P}{derivation} : \usetoken{S}{identity} सहित}

$\eq$ को !!{derivation}s में सम्मिलित करने के लिए हम केवल नए अभिगृहीत-प्रारूप
जोड़ते हैं। !!{derivation} और $\Proves$ की परिभाषा वही रहती है; अब नए
अभिगृहीतों की भी अनुमति रहती है।

\begin{defn}[!!{identity} के अभिगृहीत]
\begin{align}
\ollabel{ax:id1} & \eq[t][t], \\
\ollabel{ax:id2} & \eq[t_1][t_2] \lif (!B(t_1) \lif
  {}!B(t_2)),
\end{align}
किन्हीं बंद पदों $t$, $t_1$, $t_2$ के लिए।
\end{defn}

\begin{prop}\ollabel{prop:sound}
  अभिगृहीत \olref{ax:id1} और \olref{ax:id2} वैध हैं।
\end{prop}

\begin{proof}
  अभ्यास।
\end{proof}

\begin{prob}
\olref[fol][axd][ide]{prop:sound} सिद्ध कीजिए।
\end{prob}

\begin{prop}\ollabel{prop:iden1}
 $\Gamma \Proves \eq[t][t]$, किसी भी पद $t$ और समुच्चय~$\Gamma$ के लिए।
\end{prop}

\begin{prop}\ollabel{prop:iden2}
  यदि $\Gamma \Proves !A(t_1)$ और $\Gamma \Proves
  \eq[t_1][t_2]$, तो $\Gamma \Proves !A(t_2)$।
\end{prop}

\begin{proof}
!!{formula}
\[
(\eq[t_1][t_2] \lif (!A(t_1) \lif !A(t_2)))
\]
~\olref{ax:id2} का एक दृष्टान्त है। निष्कर्ष~\MP{} के दो प्रयोगों से मिलता है।
\end{proof}

\end{document}
